\label{chap:83-15}

\section*{CKM Mixing and Charge-Parity Violation}

The incidence sextet
\((b,p,q,h_6,o_8,t)=(4,5,7,6,8,3)\) fixes the rational transformation
\begin{equation}\label{p12-83-c15-s1:eq:constantes-ckm-matriz-racional}
 M_{\rm CKM}=
 \begin{pmatrix}
 2&-5&28\\
 0&7&-25/2\\
 1&-20&-2/3
 \end{pmatrix},
 \qquad
 \det M_{\rm CKM}=-\frac{3857}{6}\ne0.
\end{equation}
With
\[
 h_{\deg}=\frac{C^*_{\deg}}6,\qquad
 \Delta_{\deg}=\frac{180}{\pi}\Delta_4,
\]
we define
\[
 \begin{pmatrix}\theta_{12}\\\theta_{23}\\\theta_{13}\end{pmatrix}
 =M_{\rm CKM}
 \begin{pmatrix}A_{\deg}\\h_{\deg}\\\Delta_{\deg}\end{pmatrix},
 \qquad \delta_{{\rm CP},\deg}=9A_{\deg}.
\]
To prevent any ambiguity of chart, it is fixed
\[
 s_{ij}:=\sin\!\left(\frac{\pi\theta_{ij}}{180}\right),\qquad
 c_{ij}:=\cos\!\left(\frac{\pi\theta_{ij}}{180}\right).
\]
The standard parametrization produces a unitary matrix \(V\), and the
Jarlskog invariant is
\begin{equation}\label{p12-83-c15-s1:eq:constantes-jarlskog}
 J=c_{12}c_{23}c_{13}^2s_{12}s_{23}s_{13}
 \sin\!\left(\frac{\pi\delta_{{\rm CP},\deg}}{180}\right)
 =3.1189723326632974\times10^{-5}\ne0.
\end{equation}
Nonvanishing expresses CP violation in the constructed model. Comparison
with the Particle Data Group is undertaken afterward; it does not adjust
\eqref{p12-83-c15-s1:eq:constantes-ckm-matriz-racional}. The PMNS sector requires its own
mixing operator and is identified neither with CKM nor with a rank-one
diagonal phase.


\section{Dimensionless construction of the electronic scalar from APP incompatibility and the dodecaphase state}
\label{p12-83-c15-s2:chap:biografia-electron-rev6}
\index{electron!holographic construction}
\index{positron!conjugate band}
\index{mass!electronic law}

The electron appears in HMT as the point at which six
preceding and mutually typed constructions are assembled: the incompatibility between
the two arithmetic sheets of APP, the orientation induced by the involution
of the areal--volumetric hinge, two integer quotients from the APP--TPK census,
the dodecaphasic
closure of \(\alpha_{\mathrm{HMT}}\), the global defect
\(\Delta_4\), and the unique center of the nonadic atlas. The route band then adds
electron--positron conjugation. The MeV chart and the comparison
with CODATA are subsequent to all of them.

This precedence is part of the mathematical content: the factors, their
types and their causal order are determined before opening the tabulated chart, which
intervenes only in the subsequent metrological comparison.

\subsection{Electronic scalar \(\mathfrak e_{\mathrm{HMT}}\) and causal graph of its factors APP–TPK, \(\alpha_{\mathrm{HMT}}\), \(\Delta_4\) and nonadic atlas}

The internal output is the dimensionless scalar.
\begin{equation}
 \label{p12-83-c15-s2:eq:electron-holografico-rev6}
 \boxed{
 \mathfrak e_{\mathrm{HMT}}
 =
 \frac{\sqrt3}{4}
 \left[
 \exp\!\left(\frac{\varphi}{\pi^2}\right)
 -\frac{396}{18}\,\alpha_{\mathrm{HMT}}^3
 \right]
 \left[
 1+\frac{270}{18}\,\Delta_4
 \right].}
\end{equation}
The identities
\[
 \frac{396}{18}=22=27-5,
 \qquad
 \frac{270}{18}=15=\binom62=3\cdot5
 \]
allow the abbreviated notation to be recovered
\[
 \mathfrak e_{\mathrm{HMT}}
 =
 \frac{\sqrt3}{4}
 \left(e^{\varphi/\pi^2}-22\alpha_{\mathrm{HMT}}^3\right)
 (1+15\Delta_4).
\]
Here the letter \(e\) inside the exponential is Euler's number; the
subscript of \(\mathfrak e_{\mathrm{HMT}}\) denotes the electronic section.

The causal graph of the equation is
\[
\begin{array}{ccl}
\text{APP projectors in dimension three}
&\longrightarrow& \sqrt3/4,\\[2mm]
\text{TPK calendar and areal--volumetric sheets}
&\longrightarrow&
F_{\mathrm{av}},\ \pi/\varphi^2,\
\varphi/\pi^2,\\[2mm]
\text{dodecaphasic closure}
&\longrightarrow&\alpha_{\mathrm{HMT}},\\[2mm]
\text{APP--TPK census and electronic family}
&\longrightarrow&396,\ 18,\ 270,\ 22,\ 15,\ 5,\\[2mm]
\text{global readings of }\pi,e\text{ and register}
&\longrightarrow&\Delta_4,\\[2mm]
\text{nonadic atlas and anti-section}
&\longrightarrow&(5,5),\
\text{electron--positron band}.
\end{array}
\]
No arrow in this diagram receives as an argument the reference electronic mass.

\subsection{1st factor: exact incompatibility of the APP leaves}

For \(d\geq2\), let
\[
 \mathscr A_d=\{1,\ldots,9\}^d
\]
with uniform measure, and let
\[
 \Sigma_d(x)=\operatorname{dr}(x_1+\cdots+x_d),
 \qquad
 \Pi_d(x)=\operatorname{dr}(x_1\cdots x_d)
\]
the additive and multiplicative observables. Their orthogonal conditional expectations
on \(\ell^2(\mathscr A_d)\) are
\(P_{\Sigma,d}\) and \(P_{\Pi,d}\). The
\cref{prop:app-no-conmutatividad} proves by exact enumeration that
\[
 \left\|[P_{\Sigma,3},P_{\Pi,3}]\right\|
 =\frac{\sqrt3}{4}.
\]

\begin{proposition}[Electronic prefactor APP content]
\label{p12-83-c15-s2:prop:prefactor-app-electron-rev6}
The prefactor \(\sqrt3/4\) of
\cref{p12-83-c15-s2:eq:electron-holografico-rev6} is the exact norm of the
noncommutativity between the sum and product sheets in dimension three.
It is not a geometric factor introduced from the electronic formula.
\end{proposition}

\begin{proof}
The joint table of \(\Sigma_3\) and \(\Pi_3\) determines the principal
angles between
\(\operatorname{Ran}P_{\Sigma,3}\) and
\(\operatorname{Ran}P_{\Pi,3}\). If \(s\) is one of its singular
values, the corresponding contribution to the commutator norm is
\(s\sqrt{1-s^2}\). Enumeration of the \(9^3\) words gives at most
\[
 s^2(1-s^2)=\frac3{16},
 \qquad s^2=\frac14.
\]
Therefore the maximum norm is \(\sqrt{3/16}=\sqrt3/4\), before
introducing \(\alpha_{\mathrm{HMT}}\), \(\Delta_4\), a unit of energy
or a reference mass.
\end{proof}

The physical function of the factor is thereby fixed: it measures the maximal incompatibility of the two arithmetic readings on the same ternary support. The electron does not inherit two independent numbers, one additive and the other multiplicative; it inherits the exact obstruction to rendering the two sheets simultaneously compatible. The prefactor preserves, within the mass law, the local memory of the APP disagreement from which the route arises.

\subsection{2nd factor: the exchange involution between \(\pi\) and \(\varphi\)}

The primitive calendar of modes
\((\Sigma,\Pi,\Pi)\) induces, on the areal and volumetric sheets, the matrix
\[
 F_{\mathrm{av}}
 =
 \begin{pmatrix}
 0&1\\
 1&1
 \end{pmatrix}.
\]
\cref{p12-83-c16-s2:thm:descenso-necesario-fav} obtains it by counting the transitions
\(\Sigma\to\Pi\), \(\Pi\to\Pi\) and \(\Pi\to\Sigma\), each once.
The matrix is not postulated from the golden ratio: its eigenvalues
\(\varphi\) and \(-\varphi^{-1}\) appear after the incidence has been
constructed.

In the second-order memory,
\(\operatorname{Sym}^2(F_{\mathrm{av}})\) has spectrum
\[
 \varphi^2,\qquad -1,\qquad\varphi^{-2}.
\]
Thus, \(\varphi^{-2}\) is the positive stable mode of the hinge. The closure
mark \(\pi\), applied just once to the areal return, produces
\[
 \frac{\pi}{\varphi^2}.
\]
This reason preserves the closure--self-scaling order. The electronic law uses its opposite orientation.

\begin{theorem}[Exact Inversion of the Electronic Hinge]
\label{p12-83-c15-s2:thm:espejo-electron-rev6}
Let
\[
 \mathscr R(x,y)=\frac{x}{y^2},
 \qquad
 \mathsf J(x,y)=(y,x).
\]
Then \(\mathsf J^2=I\) and
\[
 \mathscr R(\pi,\varphi)=\frac{\pi}{\varphi^2},
 \qquad
 (\mathscr R\circ\mathsf J)(\pi,\varphi)
 =\frac{\varphi}{\pi^2}.
\]
Consequently,
\[
 \frac{\varphi}{\pi^2}
 =
 \frac{1}{
 \varphi^3\bigl(\pi/\varphi^2\bigr)^2},
 \qquad
 e^{\varphi/\pi^2}
 =
 \exp\!\left[
 \frac{1}{
 \varphi^3\bigl(\pi/\varphi^2\bigr)^2}
 \right].
\]
\end{theorem}

\begin{proof}
The first pair of identities is obtained by applying \(\mathscr R\) before
and after the involution \(\mathsf J\). For the reparametrization,
\[
 \varphi^3
 \left(\frac{\pi}{\varphi^2}\right)^2
 =\frac{\pi^2}{\varphi},
\]
and its inverse is \(\varphi/\pi^2\). The last equality follows by applying
the exponential.
\end{proof}

The simultaneous presence of \(\pi\) and \(\varphi\) therefore has a precise structural meaning. \(\pi/\varphi^2\) reads areal closure over quadratic volumetric memory; \(\varphi/\pi^2\) reads interior self-growth over quadratic closure. They are not two similar quotients chosen for convenience. They are the two orientations of the same reader, exchanged by an exact involution. The term \(e^{\varphi/\pi^2}\) publishes the interior orientation of that hinge in the electron section.

\subsection{3rd and 4th factors: the integers 22, 15 and the center 5}

The census of the APP chart employed by the electronic family provides
three integer sums:
\[
 \Sigma(\mathrm{APP})=396,
 \qquad
 \Sigma(\mathrm{centro}\ 2\times2)=18,
 \qquad
 \Sigma(\mathrm{canal})=270.
\]
The first admits the ring decomposition.
\[
 396=18+72+108+198.
\]
The kernel \(2\times2\) is therefore an explicit substructure of the
complete chart, and the integer 270 belongs to the global channel of the same
APP--TPK census.
The sum of the even rings is
\[
 18+108=126,
 \qquad
 \frac{396}{126}=\frac{22}{7}.
\]
This last quotient is the rational closure seal of the chart, not an
identification of \(22/7\) with the number \(\pi\) generated by the coinductive
reader. It does show that the numerator \(22\) already serves an
internal function before appearing in the electronic law.

\begin{theorem}[Internal Origin of the Electronic Coefficients]
\label{p12-83-c15-s2:thm:coeficientes-electron-rev6}
In the electronic family fixed by the APP chart and its global channel,
\[
 \boxed{
 22=\frac{\Sigma(\mathrm{APP})}
 {\Sigma(\mathrm{centro})}
 =\frac{396}{18}=27-5,}
\]
\[
 \boxed{
 15=\frac{\Sigma(\mathrm{canal})}
 {\Sigma(\mathrm{centro})}
 =\frac{270}{18}=3\cdot5=\binom62.}
\]
In particular, the coefficients are determined without consulting
\(Z_0\), the electron mass or a chart of units.
\end{theorem}

\begin{proof}
The first two equalities are exact divisions of the finite census:
\[
 396=22\cdot18,
 \qquad
 270=15\cdot18.
\]
By defining
\[
 p:=27-22,
\]
one obtains \(p=5\), and then
\[
 15=3p=3\cdot5=\binom62.
\]
All quantities belong to the integer domain of the chart and the
family before metrological transduction.
\end{proof}

\subsubsection{Persistence selector in the central return}

The same pair has a dynamic construction that also fixes its sign without
consulting a mass. The four orientations of the central seed
\((5,5)\) return to that cell for the first time after twenty-seven steps. The
typed return space is
\[
 \mathcal K_e=\mathbb Z/9\mathbb Z\times\mathbb F_3,
 \qquad |\mathcal K_e|=27.
\]
Moreover, the microfiber of \(\pi\) contains the five genealogically
distinct regions
\[
 \mathscr R_\pi=3_{++}+1_{--}+1_{\perp}.
\]
The regional gauge \(\operatorname{diag}_c\) and, within a repeated phase,
the internal order of \(U_6\) define the injection
\[
 \iota:\mathscr R_\pi\hookrightarrow\mathcal K_e
\]
through the five distinct places
\[
 (4,0),\quad(5,0),\quad(6,0),\quad(6,1),\quad(6,2).
\]
The first two correspond, respectively, to the transverse region and
the mutated return; the last three separate the positive regions by their
tritic sheet. The common visible word does not intervene as the inverse of the map.

\begin{theorem}[Internal selector of the electronic coefficients]
\label{p12-83-c15-s2:thm:selector-persistencia-electron-rev7}
With the central return and the five typed regions fixed, the coefficients of
the electronic equation are
\[
 \boxed{
 n_e=-\bigl|\mathcal K_e\setminus\iota(\mathscr R_\pi)\bigr|=-22,
 \qquad
 m_e^{\rm tor}=\bigl|\mathscr R_\pi\times\mathbb F_3\bigr|=+15.}
\]
The first symbol records the boundary deficit; the second, the retained torsional memory.
\end{theorem}

\begin{proof}
The five places above are distinct, so \(\iota\) is
injective and
\[
 |\mathcal K_e\setminus\iota(\mathscr R_\pi)|=27-5=22.
\]
Each of the five regions preserves three triticale leaves, hence
\[
 |\mathscr R_\pi\times\mathbb F_3|=5\cdot3=15.
\]
The orientation rule assigns a negative sign to states removed from the
shell and a positive sign to retained memory. The construction opens only the
center, the return of twenty-seven states, and the regional identities; it does not
open an observed mass, an impedance, or a calibrated signature.
\end{proof}

The number \(22\) measures how many central kernels of weight \(18\) are recomposed by the
APP sum \(396\). The number \(15\) measures how many of those same kernels
are recomposed by the channel \(270\), and simultaneously counts the pairs of a
hexad. The number \(5\) is not an additional parameter:
\[
 5=27-22,\qquad 15=3\cdot5.
\]
The three expressions display the same closure from depth 27,
from the central normalization, and from hexadic combinatorics.

\subsubsection{Internal derivation of coefficients 22 and 15 through the APP--TPK censuses}

The historical exploration first located \(p=5\) through a finite
comparison with \(Z_0\), and from there wrote
\((27-p,3p)=(22,15)\). That procedure retains value as a history of
discovery and as an external check. It is not, however, the proof
that governs the law: the subsequent count
\[
 \frac{396}{18}=22,
 \qquad
 \frac{270}{18}=15
\]
constructs both integers within APP--TPK and recovers \(p=5\) as a
consequence. The history explains how the pair was found; the census
explains why it belongs to the equation.

\subsection{5.º factor: the cubo of the closure \(\alpha_{\mathrm{HMT}}\)}

The closure of \(\alpha_{\mathrm{HMT}}\) precedes Dimensional Mechanics.
As established in \cref{chap:alpha}, the twelve base-thousand blocks
of \(\pi\), \(e\), \(\varphi\) and the seal \(K\) combine through the
reversible carry
\[
 P_m+E_m-\Phi_m-K_m+c_{m+1}
 =A_m+1000c_m,
 \qquad c_{13}=c_1=0.
\]
The output is
\[
 A_\alpha=
 (007,297,352,569,283,800,997,285,105,472,380,663)
\]
and, by concatenation, the exact dodecaphasic window
\[
 \alpha_{12}=\pi_{12}(\alpha_{\mathrm{HMT}})
 =
 \frac{
 7297352569283800997285105472380663
 }{10^{36}}.
\]
Telescoping the carries proves the complete integer identity; the
compatible prolongation constructs
\(\alpha_{\mathrm{HMT}}=\varprojlim_N A^{(N)}\), and \(\alpha_{12}\) is its
thirty-six-digit rational projection. The
comparison with the fine-structure constant is performed subsequently and does not
intervene in its construction.

The electronic law takes the third power
\[
 \alpha_{\mathrm{HMT}}^3.
\]
Mathematically, this is the multiplicative evaluation of the third
symmetric power of the same scalar already closed; neither three
metrological copies nor three different parameters are introduced. Its contribution is
\[
 22\alpha_{\mathrm{HMT}}^3
 =
 0.00000854906599200551727014979807598228\ldots.
\]

\begin{proposition}[Statute of the Electronic Cubic Channel]
\label{p12-83-c15-s2:prop:alpha-cubica-electron-rev6}
The term \(22\alpha_{\mathrm{HMT}}^3\) of the electronic law depends
only on the prior closure of \(\alpha_{\mathrm{HMT}}\), on internal
multiplication, and on the integer \(22\) of
\cref{p12-83-c15-s2:thm:coeficientes-electron-rev6}. It does not depend on a reading of the
electron mass.
\end{proposition}

\begin{proof}
\(\alpha_{12}\) is the rational number fixed by the concatenation and carry
identity, and the compatible family of windows determines the section
\(\alpha_{\mathrm{HMT}}\). The projection \(\alpha_{12}^3\) certifies the
printed digits of \(\alpha_{\mathrm{HMT}}^3\).
The coefficient \(22\) is determined by \(396/18\). The composition of
both operations contains only data constructed before the
energy chart.
\end{proof}

This proposition fixes exactly what is used for the electron. The third power is the cubic degree of this channel in the published law. There is no need here to postulate a universal rule according to which every occurrence of \(\alpha^d\) automatically represents a dimension \(d\); a universal statement of that kind would require its own theorem of multiplicativity among degrees. The electron genealogy is complete without turning that generalization into a premise.

\subsection{6th factor: minimal torsion and global defect \(\Delta_4\)}
\index{minimum torsion!\(\Delta_4\)}

The terms \emph{minimal torsion} and \emph{global defect}
denote a single invariant here, not two distinct corrections. The object
used by the mass law is
\begin{equation}
 \label{p12-83-c15-s2:eq:delta4-electron-rev6}
 \boxed{
 \Delta_4
 =
 \frac{\pi-e\log\pi}{270}
 \left(1+\frac{\pi}{729}\right).}
\end{equation}
Its inputs are the already fixed HMT readings of \(\pi\) and \(e\), the global
channel \(270\), and the interior cell \(729\). Consequently,
\[
 \Delta_4
 =
 0.00011119642269214005405113478612067909\ldots.
\]
The first ratio measures the signed defect \(\pi-e\log\pi\) in units of the 270 channel; the second factor couples that ratio to the interior cell through \(\pi/729\). This reading adds no parameters to the formula: it separates its two constituent operations. The \cref{p3-20260826:chap:medida-accion-positiva} shows that it is a pre-dimensional mathematical input and distinguishes it from the operator that eventually realizes it as a change of sheet.

\begin{proposition}[Globality of the Electronic Defect]
\label{p12-83-c15-s2:prop:delta4-global-electron-rev6}
\(\Delta_4\) is a global invariant preceding the species. In the electronic
law it enters through
\[
 1+15\Delta_4
 =
 1.00166794634038210081076702179181019\ldots.
\]
It is not a corrector chosen after knowing the residue of the electron mass.
\end{proposition}

\begin{proof}
The \cref{p12-83-c15-s2:eq:delta4-electron-rev6} contains neither a species label, a
mass, nor a unit. The integer \(15\) has already been constructed by \(270/18\).
Therefore, the composition \(1+15\Delta_4\) is determined by the
global data and the internal census.
\end{proof}

\(\Delta_4\) must be distinguished from the other blocks of the mass character and from the other notions of torsion. In particular, it is neither \(s_{270}\), nor the quadratic invariant \(135\Delta_4^2\), nor the angular residue \(\varepsilon_{\rm res}\), nor the geometric Cartan--Holst torsion. Substituting one for another would repeat or alter the level of the correction. For the electron, \(\Delta_4\) preserves the global memory of refinement; the factor \(15\) states how many times the electron family reads it. The qualifier \emph{minimal} preserves the historical name of this invariant and is not used as a substitute for a variational theorem of minimality.

\subsection{Uniqueness of the central cell (5,5) as the genealogical support of the electron}

Let
\[
 \mathcal C_{81}=\{1,\ldots,9\}^2
\]
the nonadic atlas and
\[
 \rho(r,c)=(10-r,10-c)
\]
its central inversion. \cref{thm:censo-atlas-81} proves that the atlas
contains a unique generation-zero cell and that this cell bears the positional
label \(e_{\mathrm{core}}\).

\begin{proposition}[Uniqueness of the Electronic Center]
\label{p12-83-c15-s2:prop:centro-electron-rev6}
The involution \(\rho\) has a unique fixed point:
\[
 \boxed{(r,c)=(5,5).}
\]
Its uniqueness is combinatorial and does not depend on the value of
\(\mathfrak e_{\mathrm{HMT}}\).
\end{proposition}

\begin{proof}
The equation \((r,c)=\rho(r,c)\) is equivalent to
\[
 2r=10,\qquad2c=10.
\]
In \(\{1,\ldots,9\}^2\) it has the unique solution \(r=c=5\). The evaluation
of the mass does not appear in the argument.
\end{proof}

\begin{corollary}[Universality of the Central Section]
\label{p12-83-c15-s2:cor:universalidad-electron-central-rev6}
Every admissible automorphism of the atlas that intertwines the central inversion
preserves the cell \((5,5)\).
\end{corollary}

\begin{proof}
If \(f\rho=\rho f\) and \(\rho(5,5)=(5,5)\), then
\(\rho f(5,5)=f\rho(5,5)=f(5,5)\). Therefore \(f(5,5)\) is another fixed point
of \(\rho\); uniqueness forces \(f(5,5)=(5,5)\).
\end{proof}

This is the structural reason why the electronic class is the same in
every realization preserving the atlas and its inversion: a local decimal does not
recur by chance; rather, every admissible chart must
transport the same fixed point and the same chain of invariants. The
physical reading of this uniqueness identifies each electron with a realization of that
central class; the proposition proves the internal universality of the atlas.

In the electronic family fixed by the construction, this point is the geometric
anchor of the section. Four objects should not be confused:
\[
 (5,5),\qquad
 e_{\mathrm{core}},\qquad
 v_e=(0,0,0,0,0),\qquad
 \text{the raw signature of the central cell}.
\]
The first is a position; the second, its internal label; the third,
the chosen origin in the fine action lattice; the fourth preserves the
raw atlas data and need not be the zero vector. Their
coordination explains why the electron occupies the center without erasing the
genealogy of the representations describing it.

\paragraph{Oriented electronic routes \(e_{++},e_{--}\), envelope \(e_{\mathrm{env}}\), action origin and signatures of the central cell}
The central cell carries two oriented histories of twelve events. Writing
\(v^2\) for the concatenation of two copies of a sextuple, the active
route of orientation \({++}\), its opposite history, and the historical
envelope are, respectively,
\[
\begin{aligned}
 e_{\rm act}=e_{++}
   &=(2,2,5,-4,-3,-2)^2,\\
 e_{--}
   &=(4,3,2,-2,-2,-5)^2,\\
 e_{\rm env}
   &=(4,3,5,-4,-3,-5)^2.
\end{aligned}
\tag{E-centro-1}\label{p12-83-c15-s2:eq:electron-rutas-finas-rev8}
\]
The two stories publish the same word
\[
 \tau_{12}=\texttt{+++---+++---},
\]
but they are not the same route. The envelope takes the absolute maximum
coordinate by coordinate while preserving the sign; it records amplitudes, but
flattens the genealogical order and does not replace \(e_{\rm act}\).

The other readers of the center live in distinct codomains:
\[
\begin{array}{rcll}
 v_e&=&(0,0,0,0,0),
   &\text{affine origin of the action lattice},\\
 R(5,5)&=&(24,0,12,0,-12),
   &\text{raw APP signature of the cell},\\
 \operatorname{Sig}^{\Gamma_{108}}_{\rm masa}
   &=&(-12,-4,12,-6,12),
   &\text{even mass projection},\\
 \operatorname{Sig}^{\Gamma_{108}}_{\rm carga}
   &=&(0,0,0,0),
   &\text{odd charge projection}.
\end{array}
\tag{E-centro-2}\label{p12-83-c15-s2:eq:electron-firmas-tipadas-rev8}
\]
The fact that the charge word has zero odd signature does not turn the route into
zero, nor does it identify the raw signature with the affine origin. Position, active
route, opposite route, envelope, origin and both signatures are coordinated
objects, not six notations for the same vector.

\subsection{Electron–positron conjugate orbit on the oriented path band}

Electronic identity does not end at the central point or at the mass
value. Material conjugation is realized on dodecaphasic words. Let
\[
 C_M(\mathbf t)=-\operatorname{rev}(\mathbf t)
\]
be the oriented anti-section and let
\(\chi\in\{-1,+1\}\) be the charge orientation. The band is the class
\[
 [(\mathbf t,\chi)]_M
 =
 \{(\mathbf t,\chi),(C_M\mathbf t,-\chi)\}
\]
of \cref{def:banda-dodecafasica-ruta}. The band preserves its even
components and exhibits two sections with opposite odd components.

\begin{theorem}[Electron and positron as sections of a band]
\label{p12-83-c15-s2:thm:electron-positron-biografia-rev6}
For the band character of
\cref{thm:principio-conjugacion-masa-banda},
\[
 m(C_M\mathbf t,-\chi)=m(\mathbf t,\chi).
\]
In the electronic channel,
\[
 (n,\chi)=(-22,+1)
 \quad\longleftrightarrow\quad
 (+22,-1)
\]
and both pairs produce
\[
 \chi n=-22.
\]
Therefore, electron and positron have the same mass and opposite charge orientations.
\end{theorem}

\begin{proof}
If \(E_+\) and \(E_-\) are, respectively, the even and odd
components of the character with respect to \(C_M\), then
\[
 E_+(C_M\mathbf t)=E_+(\mathbf t),
 \qquad
 E_-(C_M\mathbf t)=-E_-(\mathbf t).
\]
From here,
\[
 E_{\mathrm{band}}(C_M\mathbf t,-\chi)
 =
 E_{\mathrm{band}}(\mathbf t,\chi),
\]
and the equality of mass is obtained by applying the exponential. In the electronic term,
\[
 (+1)(-22)=(-1)(+22)=-22.
\]
\end{proof}

The finite closure of \cref{prop:cierre-antiseccion-ct108} reinforces this
interpretation: over the eighty-one CT--108 routes, neither pure sign nor
pure reversal preserves the taxonomy, whereas
\(-\operatorname{rev}\) does preserve it, with
\[
 81=54+9+9+9.
\]
The antiparticle is therefore neither a negative mass nor an independent
copy. It is the other oriented section of the same band: it preserves the
even invariants that determine mass and reverses the odd invariants that
determine charge and orientation.

\subsection{Basal stage of the electronic closure}

\begin{theorem}[Holographic composition of the electronic section]
\label{p12-83-c15-s2:thm:cierre-electron-holografico-rev6}
With the objects constructed in the preceding sections fixed, the
dimensionless electronic section is
\[
 \mathfrak e_{\mathrm{HMT}}
 =
 \left\|[P_{\Sigma,3},P_{\Pi,3}]\right\|
 \left[
 \exp\bigl((\mathscr R\circ\mathsf J)(\pi,\varphi)\bigr)
 -\frac{\Sigma(\mathrm{APP})}
 {\Sigma(\mathrm{centro})}\,
 \alpha_{\mathrm{HMT}}^3
 \right]
 \left[
 1+
 \frac{\Sigma(\mathrm{canal})}
 {\Sigma(\mathrm{centro})}\,
 \Delta_4
 \right].
\]
Equivalently,
\[
 \boxed{
 \mathfrak e_{\mathrm{HMT}}
 =
 \frac{\sqrt3}{4}
 \left(e^{\varphi/\pi^2}
 -22\alpha_{\mathrm{HMT}}^3\right)
 (1+15\Delta_4).}
\]
Its evaluation is
\[
 \boxed{
 \mathfrak e_{\mathrm{HMT}}
 =
 0.51099892248806607250987087719134943763\ldots.}
\]
This scalar is the basal stage \(\mathfrak e_e^{(0)}\), not the governing
electronic output. The bidirectional beta closure publishes
\(0.510998950690000808608\ldots\,\mathrm{MeV}\).
\end{theorem}

\begin{proof}
The \cref{p12-83-c15-s2:prop:prefactor-app-electron-rev6} contributes
\[
 \left\|[P_{\Sigma,3},P_{\Pi,3}]\right\|=\sqrt3/4.
\]
The \cref{p12-83-c15-s2:thm:espejo-electron-rev6} contributes
\[
 (\mathscr R\circ\mathsf J)(\pi,\varphi)=\varphi/\pi^2.
\]
The \cref{p12-83-c15-s2:thm:coeficientes-electron-rev6} contributes
\[
 \frac{\Sigma(\mathrm{APP})}{\Sigma(\mathrm{centro})}=22,
 \qquad
 \frac{\Sigma(\mathrm{canal})}{\Sigma(\mathrm{centro})}=15.
\]
The dodecaphasic closure contributes the rational window \(\alpha_{12}\) of the
coinductive section \(\alpha_{\mathrm{HMT}}\), and
\cref{p12-83-c15-s2:eq:delta4-electron-rev6} contributes the global defect. Substitution
yields the abbreviated formula. Successive evaluation of
\[
\begin{aligned}
 e^{\varphi/\pi^2}
 &=1.17814494259630678081160095680888910\ldots,\\
 22\alpha_{\mathrm{HMT}}^3
 &=0.00000854906599200551727014979807598228\ldots,\\
 1+15\Delta_4
 &=1.00166794634038210081076702179181019\ldots
\end{aligned}
\]
gives the declared value.
\end{proof}

The theorem is holographic in a verifiable sense: each local factor
retains a distinct layer of the complete system. The commutator retains sheet
duality; the exponential retains the interior orientation of the
hinge \(\pi\)--\(\varphi\); \(22\) and \(15\) retain the global census and
its central kernel; \(\alpha_{\mathrm{HMT}}^3\) retains fine closure; and
\(\Delta_4\) retains global refinement. The center determines where the
section is anchored, and the band determines how it survives under conjugation.
The value does not replace this structure: it is its common evaluation.

\subsubsection{Anterior and posterior action splitting upon return}
\label{p12-83-c15-s2:sec:electron-desdoblamiento-two-way-rev9}

The regional realization of the conjugate angular coordinate \(C^*\) distinguishes two coordinates on the
same internal action line. The coordinate before the return is \(H_5\) and
the subsequent coordinate is
\[
\etaRet
 =H_5-\frac{90}{\pi}\mathscr C_\pi(\alpha_{\rm HMT}).
\]
The notation \(\etaRet\) is the canonical sexagesimal coordinate defined on the
action branch by \(\etaRet=\Cstar/2-\alphaAngle\). The preceding equality is
evaluated in that chart; it does not introduce a second return coordinate.
Therefore, the additive and multiplicative charts of the oriented splitting
are
\begin{equation}
 \boxed{
 \delta_{2w}:=H_5-\etaRet
 =\frac{90}{\pi}\mathscr C_\pi(\alpha_{\rm HMT}),
 \qquad
 R_{\rm act}:=\frac{\hbar^{\rm pre}_{\rm HMT}}
 {\hbar^{\rm ret}_{\rm HMT}}
 =\frac{H_5}{\etaRet}.}
 \label{p12-83-c15-s2:eq:electron-desdoblamiento-two-way-rev9}
\end{equation}
Both coordinates follow from a single pre-return/post-return transport.
They are not two Planck constants: the common factor
\(10^{-34}\mathcal U_S\) fixes the dimensional line and cancels exactly in
\(R_{\rm act}\). The action ellipse subsequently realizes
\(\operatorname{Area}(\theta)=\hbar_{\rm HMT}^{\rm ret}\theta\), so that
the splitting precedes the holonomy reading.

The coupling
\begin{equation}
 \kappa_e^{\beta}
 =80-54\alpha_{\rm HMT}+6\Delta_4,
 \qquad
 \mathfrak e_e^{\beta}
 =\mathfrak e_{\rm HMT}R_{\rm act}^{\kappa_e^{\beta}}
 =0.5109989506900008086\ldots
 \label{p12-83-c15-s2:eq:electron-beta-two-way-rev9}
\end{equation}
is the bidirectional quantitative closure of the electron. The two independent
readers of the electronic register constructed in the following chapter
select, without consulting the comparison mass,
\[
 K_D(\Gamma_e)=K_\Omega(\Gamma_e)=(80,54,6),
\]
and derive exactly \(\kappa_e^\beta\). The integer \(80\) comes from the
first coordinate of those readers; the word \(w_e=201101\) retains its
own residence in the channel of Euler's constant.

\begin{statusbox}[title={Internal Closure of Bidirectional Refinement}]
The identity \eqref{p12-83-c15-s2:eq:electron-desdoblamiento-two-way-rev9}, the reduction
\((80,54,6)\), and the equation
\eqref{p12-83-c15-s2:eq:electron-beta-two-way-rev9} constitute internally exact,
target--free results on the electronic record. Propagation to fibers of
higher rank is formulated by means of the multisection operators.
\end{statusbox}
